% GENERATED FILE. Edit canonical lessons, then run npm run generate:books.
% canonical-source-hash: fnv1a64:586b73e20e2b7e8f
% canonical-lessons: ZH-C452-xiaohua, ZH-C452-haoping, ZH-C452-xiaofei, ZH-C452-tianse, ZH-C452-canjinzhi, ZH-R452-review-448-452

\chapter{After the Meal}
\label{ch:zh-after-the-meal}
\hlchaptermodality{452}

\begin{chapteropening}
\textbf{By the end of this chapter:} \emph{I can talk about digesting a meal, good reviews, tips, the light of the sky and paper napkins.}

Complete the last lesson of chapter 452: I can talk about digesting a meal, good reviews, tips, the light of the sky and paper napkins.
\end{chapteropening}

\section[xiāohuà]{xiāohuà — to digest}

\label{lesson:ZH-C452-xiaohua}



\begin{quote}
Before the new one: say the Mandarin for a meatless dish, then the Mandarin for to be left over.
\end{quote}

\subsection*{You'll want to know: xiāohuà}
\emph{xiāohuà} — \textquotedblleft{}to digest\textquotedblright{}.

\begin{quote}
\emph{chī le fàn mǎshàng shuìjiào duì xiāohuà bù hǎo} — Going to bed right after a meal is bad for your digestion.
\end{quote}

Its characters come later; for now it is a word for the ear.

\begin{cousinweb}
\emph{xiāohuà} is built from \emph{xiāo} 'to melt away' and \emph{huà} 'to change', both inherited from Middle Chinese, \emph{sjew} and \emph{xwaeH} (Guangyun, in Baxter's transcription), which glosses the first 'to use up' and the second 'to change'.
\end{cousinweb}

\begin{grammarlens}[title={What you've built}]
One more word for getting around.
\end{grammarlens}

\subsection*{Your turn}
\begin{itemize}
\raggedright
  \item \emph{Say it:} \emph{xiāohuà}
  \item \emph{Say it:} \emph{xiāohuà}, once more
  \item \emph{Say it:} \emph{chī le fàn mǎshàng shuìjiào duì xiāohuà bù hǎo} — Going to bed right after a meal is bad for your digestion.
  \item \emph{Recall:} say the Mandarin for a meatless dish, then the Mandarin for to be left over, then say \emph{xiāohuà} again
\end{itemize}

\begin{tcolorbox}[breakable,colback=teal!4,colframe=teal!35!black,title={Before you move on}]
What does xiāohuà mean? (\textquotedblleft{}To digest\textquotedblright{}.) Say it once more.
\end{tcolorbox}
\section[hǎopíng]{hǎopíng — good review; favourable comment}

\label{lesson:ZH-C452-haoping}



\begin{quote}
Before the new one: say the Mandarin for to be left over, then the Mandarin for to digest food.
\end{quote}

\subsection*{You'll want to know: hǎopíng}
\emph{hǎopíng} — \textquotedblleft{}good review; favourable comment\textquotedblright{}.

\begin{quote}
\emph{zhè jiā fànguǎn yǒu hěn duō hǎopíng} — This restaurant has lots of good reviews.
\end{quote}

Its characters come later; for now it is a word for the ear.

\begin{cousinweb}
\emph{hǎopíng} is built from \zh{好} 'good' and \emph{píng} 'to judge', both inherited from Middle Chinese, \emph{xawX} and \emph{bjaeng} (Guangyun, in Baxter's transcription), which glosses the second 'to weigh up'.
\end{cousinweb}

\begin{grammarlens}[title={What you've built}]
One more word for getting around.
\end{grammarlens}

\subsection*{Your turn}
\begin{itemize}
\raggedright
  \item \emph{Say it:} \emph{hǎopíng}
  \item \emph{Say it:} \emph{hǎopíng}, once more
  \item \emph{Say it:} \emph{zhè jiā fànguǎn yǒu hěn duō hǎopíng} — This restaurant has lots of good reviews.
  \item \emph{Recall:} say the Mandarin for to be left over, then the Mandarin for to digest food, then say \emph{hǎopíng} again
\end{itemize}

\begin{tcolorbox}[breakable,colback=teal!4,colframe=teal!35!black,title={Before you move on}]
What does hǎopíng mean? (\textquotedblleft{}Good review; favourable comment\textquotedblright{}.) Say it once more.
\end{tcolorbox}
\section[xiǎofèi]{xiǎofèi — tip; a tip for the waiter}

\label{lesson:ZH-C452-xiaofei}



\begin{quote}
Before the new one: say the Mandarin for to digest food, then the Mandarin for a good review online.
\end{quote}

\subsection*{You'll want to know: xiǎofèi}
\emph{xiǎofèi} — \textquotedblleft{}tip; a tip for the waiter\textquotedblright{}.

\begin{quote}
\emph{zài Zhōngguó chīfàn yìbān bù gěi xiǎofèi} — In China you don't usually leave a tip when you eat out.
\end{quote}

Its characters come later; for now it is a word for the ear.

\begin{cousinweb}
\emph{xiǎofèi} is built from \zh{小} 'small' and \emph{fèi} 'a cost', both inherited from Middle Chinese, \emph{sjewX} and \emph{phj+jH} (Guangyun, in Baxter's transcription), which glosses the second 'to spend'.
\end{cousinweb}

\begin{grammarlens}[title={What you've built}]
One more word for getting around.
\end{grammarlens}

\subsection*{Your turn}
\begin{itemize}
\raggedright
  \item \emph{Say it:} \emph{xiǎofèi}
  \item \emph{Say it:} \emph{xiǎofèi}, once more
  \item \emph{Say it:} \emph{zài Zhōngguó chīfàn yìbān bù gěi xiǎofèi} — In China you don't usually leave a tip when you eat out.
  \item \emph{Recall:} say the Mandarin for to digest food, then the Mandarin for a good review online, then say \emph{xiǎofèi} again
\end{itemize}

\begin{tcolorbox}[breakable,colback=teal!4,colframe=teal!35!black,title={Before you move on}]
What does xiǎofèi mean? (\textquotedblleft{}Tip; a tip for the waiter\textquotedblright{}.) Say it once more.
\end{tcolorbox}
\section[\texorpdfstring{\zh{天色} (tiānsè)}{天色 (tiānsè)}]{\zh{天色} (tiānsè) — the light of the sky; how dark the sky is getting}

\label{lesson:ZH-C452-tianse}



\begin{quote}
Before the new one: say the Mandarin for a good review online, then the Mandarin for a tip for the waiter.
\end{quote}

\subsection*{You'll want to know: \zh{天色}}
\textbf{\zh{天色}} — \emph{tiānsè} — \textquotedblleft{}the light of the sky; how dark the sky is getting\textquotedblright{}.

\begin{quote}
\textbf{\zh{天色不早了我们坐车回酒店}} \emph{tiānsè bù zǎo le wǒmen zuò chē huí jiǔdiàn} — It's getting late, so we're heading back to the hotel.
\end{quote}

\begin{cousinweb}
\textbf{\zh{天色}} (\emph{tiānsè}) is built from \zh{天} 'the sky' and \zh{色} 'a colour', both inherited from Middle Chinese, \emph{then} and \emph{srik} (Guangyun, in Baxter's transcription).
\end{cousinweb}

\begin{grammarlens}[title={What you've built}]
One more word for getting around.
\end{grammarlens}

\subsection*{Your turn}
\begin{itemize}
\raggedright
  \item \emph{Say it:} \emph{tiānsè}
  \item \emph{Say it:} \emph{tiānsè}, once more
  \item \emph{Say it:} \emph{tiānsè bù zǎo le wǒmen zuò chē huí jiǔdiàn} — It's getting late, so we're heading back to the hotel.
  \item \emph{Recall:} say the Mandarin for a good review online, then the Mandarin for a tip for the waiter, then say \emph{tiānsè} again
\end{itemize}

\begin{tcolorbox}[breakable,colback=teal!4,colframe=teal!35!black,title={Before you move on}]
What does \zh{天色} mean? (\textquotedblleft{}The light of the sky; how dark the sky is getting\textquotedblright{}.) Say it once more.
\end{tcolorbox}
\section[cānjīnzhǐ]{cānjīnzhǐ — paper napkin}

\label{lesson:ZH-C452-canjinzhi}



\begin{quote}
Before the new one: say the Mandarin for a tip for the waiter, then the Mandarin for how dark the sky is getting.
\end{quote}

\subsection*{You'll want to know: cānjīnzhǐ}
\emph{cānjīnzhǐ} — \textquotedblleft{}paper napkin\textquotedblright{}.

\begin{quote}
\emph{fúwùyuán qǐng gěi wǒ jǐ zhāng cānjīnzhǐ} — Excuse me, could I have a few napkins, please?
\end{quote}

Its characters come later; for now it is a word for the ear.

\begin{cousinweb}
\emph{cānjīnzhǐ} is built from \emph{cān} 'a meal', \emph{jīn} 'a cloth' and \emph{zhǐ} 'paper', all inherited from Middle Chinese, \emph{tshan}, \emph{kin} and \emph{tsyeX} (Guangyun, in Baxter's transcription), which glosses the first 'to swallow'.
\end{cousinweb}

\begin{grammarlens}[title={What you've built}]
One more word for getting around.
\end{grammarlens}

\subsection*{Your turn}
\begin{itemize}
\raggedright
  \item \emph{Say it:} \emph{cānjīnzhǐ}
  \item \emph{Say it:} \emph{cānjīnzhǐ}, once more
  \item \emph{Say it:} \emph{fúwùyuán qǐng gěi wǒ jǐ zhāng cānjīnzhǐ} — Excuse me, could I have a few napkins, please?
  \item \emph{Recall:} say the Mandarin for a tip for the waiter, then the Mandarin for how dark the sky is getting, then say \emph{cānjīnzhǐ} again
\end{itemize}

\begin{tcolorbox}[breakable,colback=teal!4,colframe=teal!35!black,title={Before you move on}]
What does cānjīnzhǐ mean? (\textquotedblleft{}Paper napkin\textquotedblright{}.) Say it once more.
\end{tcolorbox}
\section[(dialogue)]{Looking back: from At the Market to After the Meal}

\label{lesson:ZH-R452-review-448-452}



\begin{quote}
Say the Mandarin for a market stall and a marketplace.
\end{quote}

\subsection*{Your turn}
\begin{itemize}
\raggedright
  \item \emph{Say it:} a market stall, a marketplace, the built-up part of a city, a market held at dawn, to set up a street stall
  \item \emph{Say it:} a night market, one's taste in food, a distinctive local flavour, light and mild (of food), greasy (of food)
  \item \emph{Say it:} seafood, a distinguishing trait, specially for that reason, authentic, a restaurant's signature dish
  \item \emph{Say it:} to eat no meat, a special character of its own, outstanding, a meatless dish, to be left over
  \item \emph{Say it:} to digest food, a good review online, a tip for the waiter, how dark the sky is getting, a paper napkin
  \item \emph{Read it:} \zh{市} \zh{味} \zh{特} \zh{色} — each one aloud, then a word that holds it
\end{itemize}

\begin{tcolorbox}[breakable,colback=teal!4,colframe=teal!35!black,title={Before you move on}]
A market stall? (\emph{tānwèi}.) A marketplace? (\textbf{\zh{市场}}, \emph{shìchǎng}.) The built-up part of a city? (\textbf{\zh{市区}}, \emph{shìqū}.) A market held at dawn? (\textbf{\zh{早市}}, \emph{zǎoshì}.) To set up a street stall? (\emph{bǎitān}.) A night market? (\emph{yèshì}.) One's taste in food? (\textbf{\zh{口味}}, \emph{kǒuwèi}.) A distinctive local flavour? (\textbf{\zh{风味}}, \emph{fēngwèi}.) Light and mild (of food)? (\emph{qīngdàn}.) Greasy (of food)? (\emph{yóunì}.) Seafood? (\emph{hǎixiān}.) A distinguishing trait? (\textbf{\zh{特点}}, \emph{tèdiǎn}.) Specially for that reason? (\textbf{\zh{特意}}, \emph{tèyì}.) Authentic? (\emph{zhèngzōng}.) A restaurant's signature dish? (\emph{zhāopáicài}.) To eat no meat? (\emph{chīsù}.) A special character of its own? (\textbf{\zh{特色}}, \emph{tèsè}.) Outstanding? (\textbf{\zh{出色}}, \emph{chūsè}.) A meatless dish? (\emph{sùcài}.) To be left over? (\emph{shèng}.) To digest food? (\emph{xiāohuà}.) A good review online? (\emph{hǎopíng}.) A tip for the waiter? (\emph{xiǎofèi}.) How dark the sky is getting? (\textbf{\zh{天色}}, \emph{tiānsè}.) A paper napkin? (\emph{cānjīnzhǐ}.)
\end{tcolorbox}
